\chapter{Objetivos de Desarrollo Sostenible (ODS) y Aspectos Sociales}
«Contextualizar los ODS que "tocan" a este trabajo . Extensión máxima: 1 página.»\\

\textbf{Contenido de la sección:}
\begin{itemize}
    \item Objetivos de desarrollo sostenible de la Agenda 2023
    \begin{itemize}
        \item \href{https://www.un.org/sustainabledevelopment/es/sustainable-development-goals/}{Link: Objetivos ODS}
    \end{itemize}
    \item Impacto social y medioambiental
\end{itemize}

Reflexión sobre los impactos sociales y ambientales del proyecto realizado, así como sobre los aspectos relativos a la responsabilidad ética y profesional que pudieran estar relacionados con el mismo

\begin{IAtext}

\section{Objetivos de Desarrollo Sostenible relacionados}

\begin{itemize}
    \item \textbf{ODS 3 — Salud y bienestar}: es el objetivo más directamente relacionado con este trabajo. Un dispositivo portátil y de bajo coste para el análisis dinámico de la pisada facilita la detección temprana de alteraciones en el patrón de marcha asociadas a patologías como el pie diabético o el Parkinson, contribuyendo potencialmente a un diagnóstico más temprano y a un seguimiento más frecuente y accesible de pacientes en tratamiento o rehabilitación, fuera del entorno estrictamente clínico.
    \item \textbf{ODS 9 — Industria, innovación e infraestructura}: el proyecto emplea fabricación aditiva (impresión 3D) y componentes electrónicos de bajo coste y propósito general para construir un dispositivo de sensorización que tradicionalmente ha requerido infraestructura de laboratorio especializada, contribuyendo a la innovación en el diseño de dispositivos médicos y deportivos accesibles.
    \item \textbf{ODS 10 — Reducción de las desigualdades}: al reducir el coste de materiales del dispositivo y eliminar la dependencia de equipamiento ortopédico adicional, el sistema desarrollado amplía potencialmente el acceso al análisis de la pisada a centros de rehabilitación, consultas de fisioterapia o deportistas amateur con recursos limitados, que de otro modo quedarían excluidos del uso de plataformas de laboratorio o soluciones comerciales de gama alta.
\end{itemize}

\section{Impacto social y medioambiental}

Desde el punto de vista social, el sistema desarrollado persigue democratizar el acceso al análisis dinámico de la pisada, llevándolo más allá del entorno clínico especializado hacia contextos de atención primaria, rehabilitación y práctica deportiva no profesional. La incorporación de persistencia local de datos en la propia aplicación, independiente de la disponibilidad de la plataforma en la nube, refuerza además la autonomía del usuario final sobre su propia información de salud, permitiéndole consultar su historial sin depender de un tercero.

Desde el punto de vista medioambiental, el uso de impresión 3D para la fabricación de la plantilla y la tobillera reduce el desperdicio de material asociado a procesos de fabricación industrial tradicionales orientados a grandes series, si bien introduce la necesidad de valorar el ciclo de vida y la reciclabilidad de los filamentos plásticos empleados. El bajo consumo energético perseguido en la optimización del firmware, además de prolongar la autonomía de la batería del dispositivo, reduce indirectamente la frecuencia de recarga y, con ella, el impacto energético asociado al uso continuado del dispositivo.

\section{Reflexión ética y profesional}

El desarrollo de dispositivos de monitorización de la salud conlleva una responsabilidad ética específica en torno al tratamiento de los datos personales recogidos (en este caso, datos biomecánicos y de actividad física, potencialmente sensibles si se emplean con fines de diagnóstico clínico). Aunque el alcance de este Trabajo de Fin de Grado es de naturaleza técnica y no incluye un despliegue clínico real, cualquier evolución del sistema hacia un uso en pacientes reales debería incorporar mecanismos explícitos de consentimiento informado, cifrado de las comunicaciones y cumplimiento de la normativa de protección de datos aplicable (en el ámbito europeo, el Reglamento General de Protección de Datos), aspectos que quedan fuera del alcance funcional actual del prototipo pero que resulta pertinente señalar como responsabilidad profesional de cara a un eventual desarrollo futuro del proyecto.

\end{IAtext}